\chapter{The interaction term}\label{ch:cI3}
\chapterbrief
{an evolution factor $Q=\tau-a(\xi)$, $a=\sum_{n=1}^Na_n$ with $a_n\in\C[\xi]$ the terms; a bounded rectangle $\Omega=I\times J$; $u\in\ker Q$, base time $t_0$, $u_0=u(t_0,\cdot)$, $s=t-t_0$.}
{the interaction term compares the full solution with the sum of the single-term evolutions. We prove its exact form and cost in a class of exponential type, that it is not continuous in Sobolev norms and cannot be regularised, that the Dyson series converges, and that non-commuting terms with an infinite adjoint orbit admit no finite product formula.}
{Dim.~A19--A24. They prove Chapter~8 of the paper.}

\begin{notebox}{The class $\Bclass$ and two lemmas used throughout}
For $\rho>0$,
\[
\|f\|_{B_\rho}=\sup_{j\ge0}\rho^{-j}\|f^{(j)}\|_{L^2(J)},\qquad B_\rho(J)=\{f:\ \|f\|_{B_\rho}<\infty\},
\]
and $\Bclass_\rho(\Omega)$ likewise with $\sup_{i,j}\rho^{-i-j}\|\partial_t^i\partial_x^ju\|_{L^2}$; $\Bclass=\bigcup_\rho\Bclass_\rho$. For $A(\xi)=\sum_l\alpha_l\xi^l$ write $|A|(\rho)=\sum_l|\alpha_l|\rho^l$.

\textbf{(P1-L1)} $B_\rho$ and $\Bclass_\rho$ are Banach spaces with $\|\partial_x\|\le\rho$; $\Bclass_\rho\hookrightarrow H^m$ for every $m$; every exponential-polynomial lies in $\Bclass$; for $u\in\Bclass_\rho(\Omega)$ the trace $t\mapsto u(t,\cdot)$ is continuous into $B_{\rho'}(J)$ for every $\rho'>\rho$ \proved.

\textbf{(P1-L2)} On $B_\rho$, $\|A(\partial_x)\|\le|A|(\rho)$, and $e^{sA(\partial_x)}=\sum_k\frac{s^k}{k!}A(\partial_x)^k$ converges in operator norm. These exponentials commute, satisfy $e^{s(A+A')(\partial_x)}=e^{sA(\partial_x)}e^{sA'(\partial_x)}$, and
\[
e^{sA(\partial_x)}\big(x^le^{\zeta x}\big)=\partial_\zeta^l\big(e^{sA(\zeta)}e^{\zeta x}\big)\qquad\proved .
\]
Both are proved in Part~II (Dim.~B for P1-L1, P1-L2).
\end{notebox}

%=====================================================================
\section{Exact form and cost}

\begin{dimcard}{A19}{P1-T2}{Closed form of the interaction}{\proved}{thm:eint}
Let $u\in\Bclass_\rho(\Omega)$ with $Q(\partial)u=0$, $\rho'>\rho$, $u_n(t)=e^{s\,a_n(\partial_x)}u_0$.
\begin{enumerate}[label=(\roman*)]
\item $u_n$ is the unique solution in $\Bclass$ of $(\partial_t-a_n(\partial_x))v=0$ with $v(t_0)=u_0$, and $u(t)=e^{s\,a(\partial_x)}u_0$.
\item In $B_{\rho'}(J)$,
\[
\Eint^{\mathrm{add}}(t)=\sum_{|S|\ge2}\ \prod_{n\in S}\big(e^{s\,a_n(\partial_x)}-\Id\big)u_0,\qquad \Eint^{\mathrm{mult}}\equiv0 .
\]
\item If $u_0=\int e^{\zeta x}\dd\mu_0(\zeta)$ with $\mu_0$ compactly supported, then with $z_n=e^{s\,a_n(\zeta)}-1$,
\[
\Eint^{\mathrm{add}}(t)=\int e_{\ge2}(z)\,e^{\zeta x}\dd\mu_0(\zeta),\qquad e_{\ge2}(z)=\prod_n(1+z_n)-1-\sum_nz_n .
\]
\end{enumerate}
\end{dimcard}

\begin{proofidea}
Everything happens in a commutative Banach algebra of bounded operators. The single-term flows commute and their product is the full flow. Subtracting the ``first-order'' part $\Id+\sum_n(e^{sa_n}-\Id)$ from the expanded product leaves exactly the subsets with at least two elements.
\end{proofidea}

\begin{proof}
\step{1}{(i) Flows are unique in $\Bclass$}
By P1-L1, $u$ is a continuous curve in $B_{\rho'}$. A solution in $\Bclass$ of $\partial_tv=A(\partial_x)v$ with datum $v(t_0)$ is $e^{(t-t_0)A}v(t_0)$; uniqueness follows from the group property on $B_{\rho_2}$ for large $\rho_2$. Apply this with $A=a$ and with $A=a_n$.

\step{2}{(ii) The full flow is the product of the single flows}
By P1-L2,
\[
e^{s\,a(\partial_x)}=\prod_{n=1}^Ne^{s\,a_n(\partial_x)},
\]
so $\Eint^{\mathrm{mult}}=0$.

\step{3}{(ii) Expanding the product}
Write $e^{s\,a_n}=\Id+\Delta_n$ with $\Delta_n=e^{s\,a_n}-\Id$. The $\Delta_n$ commute, so
\[
\prod_{n=1}^N(\Id+\Delta_n)=\sum_{S\subseteq\{1,\dots,N\}}\ \prod_{n\in S}\Delta_n .
\]
The empty set gives $u_0$; the singletons give $u_n(t)-u_0$. Subtracting them from $u(t)$ leaves the sum over $|S|\ge2$.

\step{4}{(iii) Mode by mode}
$\zeta\mapsto e^{\zeta x}\in B_{\rho'}$ is continuous on the support of $\mu_0$, so the integral is a Bochner integral and bounded operators pass under it. On one mode,
\[
\prod_{n\in S}\big(e^{s\,a_n(\partial_x)}-\Id\big)e^{\zeta x}=\prod_{n\in S}z_n\;e^{\zeta x}.
\]
Summing over $|S|\ge2$ gives $e_{\ge2}(z)$.

\whyworks{commuting bounded operators obey the algebra of numbers, and the interaction is the part of a product of $N$ factors that involves at least two of them.}
\end{proof}
\uses{P1-L1, P1-L2.}

\begin{example}{A stationary state is pure interaction}
Take $a_1=-c\xi$, $a_2=D\xi^2$ and the stationary mode $u=e^{cx/D}$, so $a(c/D)=0$. With $\theta=c^2s/D$:
\[
z_1=e^{-\theta}-1,\qquad z_2=e^{\theta}-1,\qquad
\Eint^{\mathrm{add}}=z_1z_2\,u=-4\sinh^2\!\Big(\frac\theta2\Big)\,e^{cx/D}.
\]
\end{example}

%---------------------------------------------------------------------
\begin{dimcard}{A20}{P1-T3}{Cost: $O(N)$ versus $\Theta(2^N)$}{\proved}{thm:cost}
Model: $u_0=\sum_{k=1}^Mp_k(x)e^{\zeta_kx}$ with $\deg p_k\le L$; terms $a_n$ of degree $d_n$; unit-cost arithmetic and $\exp$.
\begin{enumerate}[label=(\alph*)]
\item One value of $\Eint^{\mathrm{add}}(t)$ costs
\[
O\Big(M\big[N(L+1)^2+(L+1)\textstyle\sum_n(d_n+1)\big]\Big).
\]
\item All order aggregates $E^{(k)}=\sum_{|S|=k}\prod_{n\in S}\Delta_nu_0$, $2\le k\le N$, cost $O\big(M[N^2(L+1)^2+(L+1)\sum_n(d_n+1)]\big)$.
\item The full table of cross terms has $M(2^N-N-1)$ entries and costs $\Theta\big(M\,2^N(L+1)^2\big)$.
\end{enumerate}
\end{dimcard}

\begin{proofidea}
By A19(iii), each mode contributes $e_{\ge2}(z)$ times the mode, and $e_{\ge2}$ is a product minus a sum: it can be evaluated without ever expanding the product. Polynomial prefactors are handled by truncated power series in $\zeta$.
\end{proofidea}

\begin{proof}
\step{1}{The quantity to compute for one summand}
By A19(iii) and P1-L2, the part of $\Eint$ coming from $x^le^{\zeta_kx}$ is
\[
\partial_\zeta^l\Big[e_{\ge2}\big(z(\zeta)\big)\,e^{\zeta x}\Big]_{\zeta=\zeta_k}.
\]
So we need the Taylor coefficients of $e_{\ge2}(z(\zeta))$ at $\zeta_k$ up to order $L$. Put $\eta=\zeta-\zeta_k$ and work modulo $\eta^{L+1}$.

\step{2}{(a) Four operations per summand}
\begin{enumerate}[label=(\arabic*)]
\item $a_n(\zeta_k+\eta)\bmod\eta^{L+1}$ by repeated synthetic division: $O((L+1)(d_n+1))$.
\item $\exp\big(s\,a_n\big)\bmod\eta^{L+1}$ by the recurrence $f'=g'f$: $O((L+1)^2)$.
\item $\prod_n(1+z_n)$: $N-1$ truncated products, $O(N(L+1)^2)$.
\item Subtract $1+\sum_nz_n$ and read off the coefficients.
\end{enumerate}
Summing over $n$ and over the $M$ summands gives (a).

\step{3}{(b) All orders by dynamic programming}
The aggregate of order $k$ is the elementary symmetric function $e_k(z)$, the coefficient of $\lambda^k$ in $\prod_n(1+\lambda z_n)$. Multiplying in one factor at a time, step $n$ updates at most $n$ coefficients, each a truncated product: $O(N^2(L+1)^2)$ per summand.

\step{4}{(c) The full table}
Each cross term $c_S$ is a separate output, so the output alone has $M(2^N-N-1)$ entries of size $L+1$. They are obtained by $c_S=c_{S\setminus\{\max S\}}\cdot z_{\max S}$, one truncated product per subset, which matches the lower bound up to the factor $L+1$.

\whyworks{a product can be evaluated in linear time; only listing its expansion is exponential.}
\end{proof}
\uses{A19, P1-L2.}

%=====================================================================
\section{Instability}

\begin{dimcard}{A21}{P1-P1(e)}{The interaction is not continuous in $H^m$}{\falsest{} (continuity is false; proof of the counterexample)}{thm:nocont}
\phantomsection\label{dim:P1-P1e}%
Let $Q=\tau+c\xi-D\xi^2$ with $c,D>0$, terms $a_1=-c\xi$ and $a_2=D\xi^2$, $I=(0,T)$, and any base time $t_0\in\bar I$. For every $m\in\N$,
\[
\sup\Big\{\frac{\|\Eint^{\mathrm{add}}\|_{L^2(\Omega)}}{\|u\|_{H^m(\Omega)}}:\ u\in\EP,\ Qu=0,\ u\ne0\Big\}=\infty .
\]
Hence $u\mapsto\Eint^{\mathrm{add}}$ has no continuous extension from $\Bclass\cap\ker Q$ to $\ker_{H^m}Q$.
\end{dimcard}

\begin{proofidea}
On a bounded interval, real exponentials $e^{\zeta x}$ with large $\zeta$ are legitimate data. The diffusion flow alone amplifies them by $e^{D\zeta^2s}$, while the true solution grows only like $e^{(D\zeta^2-c\zeta)s}$. The interaction term must cancel the difference, and the ratio grows like $e^{c\zeta T}$, faster than any power of $\zeta$.
\end{proofidea}

\begin{proof}
\step{1}{The test solutions}
For real $\zeta>c/D$ let
\[
u_\zeta=e^{\alpha t+\zeta x},\qquad \alpha=D\zeta^2-c\zeta>0 .
\]
Then $Qu_\zeta=0$ and $\|u_\zeta\|_{H^m(\Omega)}\le C_m(1+\zeta)^m\,e^{\alpha T}\,\|e^{\zeta x}\|_{L^2(J)}$.

\step{2}{The interaction of one mode}
By A19(iii),
\[
\Eint^{\mathrm{add}}=e^{\alpha t_0}\,z_1z_2\,e^{\zeta x},\qquad z_1=e^{-c\zeta s}-1,\quad z_2=e^{D\zeta^2s}-1 .
\]

\step{3}{Case $t_0<T$}
Put $T'=T-t_0>0$. For $s\in[T'/2,T']$ and $\zeta$ large, $|z_1|\ge\tfrac12$ and $|z_2|\ge\tfrac12e^{D\zeta^2s}$. Integrating over $t\in[T-\zeta^{-2},T]$:
\[
\|\Eint^{\mathrm{add}}\|^2\ge c\,\zeta^{-2}\,e^{2\alpha t_0+2D\zeta^2T'}\,\|e^{\zeta x}\|^2 .
\]
Dividing by $\|u_\zeta\|_{H^m}^2$, and using $2D\zeta^2T'-2\alpha(T-t_0)=2c\zeta T'$, the ratio of squares is at least
\[
c_m\,\zeta^{-4m-2}\,e^{2c\zeta(T-t_0)}\longrightarrow\infty .
\]

\step{4}{Case $t_0=T$}
Now $s\in[-T,0]$, $z_1=e^{c\zeta|s|}-1$, and $|z_2|=1-e^{-D\zeta^2|s|}\ge\tfrac12$ for large $\zeta$. For $t\in[0,T/2]$, $|z_1|\ge\tfrac12e^{c\zeta T/2}$, and the ratio is at least $c'_m(1+\zeta)^{-4m}e^{c\zeta T}\to\infty$.

\step{5}{No continuous extension}
$u\mapsto\Eint^{\mathrm{add}}$ is linear and unbounded on a dense subspace, so it has no continuous extension.

\whyworks{a bounded interval admits real exponential modes, and on them one term alone grows much faster than the full equation.}
\end{proof}
\uses{A19.}

\begin{caution}[title={Where this does not happen}]
On $\R$ with $L^2$ data the modes are $e^{i\kappa x}$, every $|z_n|\le2$ for dissipative or dispersive terms, and the interaction is bounded: $\|\Eint\|\le2\|u_0\|$ for advection--diffusion \led{P2-P5}. The failure is a feature of \emph{bounded} domains.
\end{caution}

%---------------------------------------------------------------------
\begin{dimcard}{A22}{P2-T2}{No regularisation removes the interaction}{\proved}{thm:noreg}
\phantomsection\label{dim:P2-T4}%
\textbf{(P2-T2)} Let $\Omega=I\times J$, $a=\sum_{n=1}^Na_n$ with distinct $a_n$ and $a-a_n\ne0$ for all $n$, $D_n=\ker_{H^m}(\partial_t-a_n(\partial_x))$, $M=\overline{\sum_nD_n}$, and $u\in\ker_{H^m}(\partial_t-a(\partial_x))$. Put
\[
g=\prod_{n=1}^N(a-a_n)\in\C[\xi],\qquad K=\deg g,\qquad Q=\prod_n(\tau-a_n),\qquad k=\deg Q .
\]
\begin{enumerate}[label=(\alph*)]
\item $Q(\partial)u=g(\partial_x)u$, and $\|E^\perp(u)\|_{H^m}\ge\|g(\partial_x)u\|_{H^{m-k}}/C_Q$, where $E^\perp(u)$ is the component of $u$ orthogonal to $M$.
\item $E^\perp(u)=0\iff g(\partial_x)u=0$; this space has dimension $K$ and consists of exponential-polynomials.
\item Every split by terms, with any regulariser, leaves a residual of norm at least $\|E^\perp(u)\|$.
\end{enumerate}
\textbf{(P2-T4)} If maps $R_\alpha$ with Lipschitz constants $L_\alpha$ converge pointwise to $\Phi$ on a set $V$, then $\mathrm{Lip}(\Phi)\le\liminf L_\alpha$. Every exact selection of components has Lipschitz constant at least $\gamma_N^{-1/2}$.
\end{dimcard}

\begin{proofidea}
Every per-term piece is killed by $Q=\prod_n(\tau-a_n)$, so $Q$ kills the closed span $M$. On a solution of the full equation, $Q$ acts as the polynomial $g(\partial_x)$. So $\|Q(\partial)u\|$ measures the distance of $u$ from $M$, and only the finitely many solutions killed by $g$ have distance zero. P2-T4 is the fact that a pointwise limit of uniformly Lipschitz maps is Lipschitz with the same constant.
\end{proofidea}

\begin{proof}
\step{1}{(a) $Q$ acts as $g$ on solutions}
Divide $Q\in\C[\xi][\tau]$ by the monic $\tau-a(\xi)$:
\[
Q=(\tau-a)H+r(\xi),\qquad r(\xi)=Q\big(a(\xi),\xi\big)=\prod_n(a-a_n)(\xi)=g(\xi).
\]
Since constant-coefficient operators commute in $\mathcal D'$, $Q(\partial)u-g(\partial_x)u=H(\partial)(\partial_t-a(\partial_x))u=0$.

\step{2}{(a) The distance from $M$}
Each $v\in D_n$ is killed by its own factor $\tau-a_n$ of $Q$, hence by $Q$; by continuity $Q$ kills $M$. Write $u=P_Mu+E^\perp(u)$. The map $Q(\partial):H^m\to H^{m-k}$ is bounded with norm $C_Q$ (\led{P2-L1}), so
\[
\|g(\partial_x)u\|_{H^{m-k}}=\|Q(\partial)E^\perp(u)\|_{H^{m-k}}\le C_Q\,\|E^\perp(u)\|_{H^m}.
\]

\step{3}{(b) The exceptional space}
If $E^\perp(u)=0$, Step~2 gives $g(\partial_x)u=0$, so $u$ is killed by the ideal $\mathfrak J=(\tau-a,\ g)$. The map $\ring\to\C[\xi]$, $\tau\mapsto a(\xi)$, is onto with kernel $(\tau-a)$, so
\[
\ring/\mathfrak J\cong\C[\xi]/(g),\qquad \dimc=K .
\]
By A12, the solutions form a $K$-dimensional space of exponential-polynomials. Conversely, such $u$ split by factors of $Q$ (A13) with components in $\Bclass\subset H^m$ (\led{P1-P2(c)}), so $u\in M$.

\step{4}{(c) Any regulariser}
For any representation $S\omega$, $u-S\omega\in u+M$, so $\|u-S\omega\|\ge\mathrm{dist}(u,M)=\|E^\perp(u)\|$.

\step{5}{(P2-T4) Lipschitz constants pass to the limit}
For $v,w\in V$: $\|R_\alpha v-R_\alpha w\|\le L_\alpha\|v-w\|$, and the left side tends to $\|\Phi v-\Phi w\|$. Taking the supremum over $v\ne w$ gives $\mathrm{Lip}(\Phi)\le\liminf_\alpha L_\alpha$. For an exact selection $\sigma$ ($S\sigma(u)=u$), $\sigma(u)-\sigma(0)$ is a preimage of $u$, so its norm is at least that of the minimal preimage, $\|S^\dagger u\|$; hence $\mathrm{Lip}(\sigma)\ge\|S^\dagger\|=\gamma_N^{-1/2}$ (A17).

\whyworks{the per-term pieces all satisfy one extra equation, and on solutions of the full equation that extra equation is a nonzero polynomial in $\partial_x$.}
\end{proof}
\uses{A12, A13, A17, P2-L1, P1-P2.}

\begin{example}{Advection--diffusion}
For $a_1=-c\xi$ and $a_2=D\xi^2$: $g=(D\xi^2)(-c\xi)=-cD\xi^3$, $K=3$, and $Q(\partial)u=-cD\,u_{xxx}$. The exceptional space is
\[
\{E^\perp=0\}=\spn\big\{1,\ x-ct,\ (x-ct)^2+2Dt\big\},
\]
with, for example,
\[
(x-ct)^2+2Dt=\underbrace{\Big[(x-ct)^2-\tfrac{2D}{c}(x-ct)\Big]}_{\in D_1}+\underbrace{\tfrac{2D}{c}\,x}_{\in D_2}.
\]
\end{example}

%=====================================================================
\section{Perturbation and non-commuting terms}

\begin{dimcard}{A23}{P3-T1}{Dyson--Phillips converges for every time}{\proved}{thm:dyson}
\textbf{(a)} On $X=L^2(\R)$ let $A_0=a(\partial_x)$, $a(\xi)=\sum_{l\le q}\alpha_l\xi^l$, $\alpha_q\ne0$, and $V=\sum_{k\le r}b_k(x)\partial_x^k$ with bounded $b_k$ and $r<q$.
\begin{enumerate}[label=(a\arabic*)]
\item The condition (Par): $\mathrm{Re}\,a(i\kappa)\le\omega-\delta|\kappa|^q$ for some $\delta>0$, holds iff $q$ is even and $\mathrm{Re}(\alpha_qi^q)<0$.
\item Under (Par), the Dyson--Phillips series of $e^{t(A_0+\varepsilon V)}$ converges in operator norm, uniformly on every $[0,T]$, is entire in $\varepsilon$, and its sum is the semigroup generated by $A_0+\varepsilon V$.
\end{enumerate}
\textbf{(b)} The same holds on $L^2(0,L)$ for $A_0=D\partial_x^2$ with Dirichlet conditions and $V=b_1\partial_x+b_0$.
\end{dimcard}

\begin{proofidea}
The principal term smooths: $\|\partial^kS_0(s)\|$ blows up only like $s^{-k/q}$. Since the perturbation has order $r<q$, the singularity $s^{-r/q}$ is integrable. Iterating gives Beta integrals, and the series is bounded by a Mittag-Leffler function, which is entire.
\end{proofidea}

\begin{proof}
\step{1}{The abstract estimate (P3-L2)}
Suppose $\|S_0(t)\|\le Me^{\omega t}$ and, for $s\in(0,T]$,
\[
\|VS_0(s)\|\le K\,s^{-\beta}e^{\omega s},\qquad \beta\in[0,1),\quad \gamma=1-\beta .
\]
Define $S^{(0)}=S_0$ and $S^{(n)}(t)x=\int_0^tS^{(n-1)}(t-s)\,VS_0(s)x\dd s$. By induction,
\[
\|S^{(n)}(t)\|\le Me^{\omega t}\,\frac{\big(K\Gamma(\gamma)\big)^n\,t^{n\gamma}}{\Gamma(1+n\gamma)} .
\]
The induction step is the Beta integral
\[
\int_0^t(t-s)^{(n-1)\gamma}s^{\gamma-1}\dd s=t^{n\gamma}\,\frac{\Gamma\big((n-1)\gamma+1\big)\Gamma(\gamma)}{\Gamma(n\gamma+1)} .
\]

\step{2}{Convergence and analyticity}
The bound sums to $Me^{\omega t}E_\gamma\big(K\Gamma(\gamma)t^\gamma\big)$, where $E_\gamma(w)=\sum_nw^n/\Gamma(1+n\gamma)$ is the Mittag-Leffler function, entire by Stirling. Replacing $V$ by $\varepsilon V$ gives a power series in $\varepsilon$ with infinite radius.

\step{3}{The sum is the right semigroup}
By Miyadera--Voigt (Engel--Nagel, Thm.~III.3.14, Cor.~III.3.15--16) \cf, $A_0+\varepsilon V$ generates a semigroup $S$ satisfying the variation-of-constants formula. The series satisfies the same formula; the difference $\Delta$ satisfies $\Delta(t)=\int_0^t\Delta(t-s)\,\varepsilon VS_0(s)\dd s$, and iterating with the Beta integral gives $\|\Delta\|\le C\,(|\varepsilon|K\Gamma(\gamma))^nt^{n\gamma}/\Gamma(1+n\gamma)\to0$.

\step{4}{(a1) The symbol criterion}
Let $p(\kappa)=\mathrm{Re}\,a(i\kappa)=\mathrm{Re}(\alpha_qi^q)\kappa^q+p_{<q}(\kappa)$. If $q$ is even and $\mathrm{Re}(\alpha_qi^q)=-2\delta<0$, then $p(\kappa)\le-2\delta\kappa^q+C(1+|\kappa|^{q-1})\le\omega-\delta|\kappa|^q$. Conversely, (Par) forces $p\to-\infty$ like $-\delta|\kappa|^q$ on both sides, so $\deg p=q$, the leading coefficient is negative, and $q$ cannot be odd.

\step{5}{(a2) The hypothesis of Step~1 holds with $\beta=r/q$}
By Plancherel, for $s>0$,
\[
\|\partial^kS_0(s)\|=\sup_\kappa|\kappa|^ke^{s\,p(\kappa)}\le e^{\omega s}\sup_{y\ge0}y^ke^{-s\delta y^q}=e^{\omega s}\Big(\frac{k}{q\delta e\,s}\Big)^{k/q}.
\]
With $\|b_k\cdot\|\le\|b_k\|_\infty$ and $s^{-k/q}\le T^{(r-k)/q}s^{-r/q}$ on $(0,T]$, this gives the bound of Step~1 with $\beta=r/q<1$.

\step{6}{(b) Dirichlet}
In the sine basis $e_n=\sqrt{2/L}\sin(n\pi x/L)$,
\[
\|\partial S_0(s)f\|^2=\sum_n|f_n|^2\mu_n^2e^{-2D\mu_n^2s}\le\frac{\|f\|^2}{2eDs},
\]
so Step~1 applies with $\beta=\tfrac12$.

\whyworks{a smoothing principal term pays for a lower-order perturbation, and the payment is an integrable singularity in time.}
\end{proof}
\uses{Engel--Nagel (Miyadera--Voigt).}

\begin{caution}[title={The hypotheses cannot be dropped \led{P3-C1}}]
With a perturbation of the same order ($A_0=\partial^2$, $V=-2\partial^2$) the terms have norm $(2n)^ne^{-n}/n!\sim2^n/\sqrt{2\pi n}$ and the series diverges. With a dispersive principal term ($A_0=i\partial^2$, $V=i\partial$) the sum does not even generate a semigroup on $L^2$.
\end{caution}

%---------------------------------------------------------------------
\begin{dimcard}{A24}{P1-T5}{No finite product of flows for an infinite adjoint orbit}{\proved}{thm:lie}
Let $\mathcal A$ be a unital associative $\C$-algebra and $X,Y\in\mathcal A$ such that
\[
X,\quad Y_0=Y,\quad Y_1=\mathrm{ad}_XY,\quad Y_2=\mathrm{ad}_X^2Y,\ \dots
\]
are linearly independent. Then there is no finite word $(Z_1,\dots,Z_d)\in\{X,Y\}^d$, no common polydisc $U=\{|t|<r_0\}\times\{|\varepsilon|<\varepsilon_0\}$, and no functions $g_i$ holomorphic on $U$ with $g_i(0,\varepsilon)=0$ and $g_i(t,0)\equiv0$ whenever $Z_i=Y$, such that
\[
e^{t(X+\varepsilon Y)}=\prod_{i=1}^de^{g_i(t,\varepsilon)Z_i}\qquad\text{in }\mathcal A[[t]]\ \text{for all }|\varepsilon|<\varepsilon_0 .
\]
\end{dimcard}

\begin{proofidea}
Differentiate the identity in $\varepsilon$ at $\varepsilon=0$. The left side produces the moments $t^{k+1}/(k+1)$ of the independent elements $Y_k$. The right side, being a finite word, produces a sum of at most $p$ exponential-type terms, whose Hankel matrices have rank at most $p$. But the Hankel matrices of the moments are Hilbert matrices, of full rank.
\end{proofidea}

\begin{proof}
\step{1}{At zero coupling}
For $\varepsilon=0$ the $Y$-factors are the identity, so $e^{tX}=e^{G(t)X}$ with $G=\sum_{Z_i=X}g_i(t,0)$. If $G\ne t$, let $m$ be the order of $G-t$; the coefficient of $t^m$ in $e^{(G-t)X}=1$ gives $(G-t)_mX=0$, impossible since $X\ne0$. So $G=t$.

\step{2}{Derivative in $\varepsilon$: the left side}
Coefficients of $t^m$ on both sides live in the finite-dimensional span of words of length $\le m$, and depend holomorphically on $\varepsilon$, so we may differentiate. By the formal Duhamel formula,
\[
\partial_\varepsilon e^{t(X+\varepsilon Y)}\Big|_{\varepsilon=0}=e^{tX}\int_0^te^{-\sigma X}Ye^{\sigma X}\dd\sigma=e^{tX}\sum_{k\ge0}(-1)^k\frac{t^{k+1}}{(k+1)!}\,Y_k .
\]

\step{3}{Derivative in $\varepsilon$: the right side}
At $\varepsilon=0$ all factors are exponentials of $X$. The $X$-factors contribute $r(t)Xe^{tX}$ with $r=\sum_{Z_i=X}\partial_\varepsilon g_i(t,0)$. The $Y$-factors contribute $q_i(t)\,e^{G_iX}Ye^{(t-G_i)X}$, with $q_i=\partial_\varepsilon g_i(t,0)$ and $G_i$ the sum of the earlier $X$-times. Multiply on the left by $e^{-tX}$ and use
\[
e^{-\tau X}\,Y\,e^{\tau X}=\sum_k(-1)^k\frac{\tau^k}{k!}\,Y_k,\qquad \tau_i=t-G_i .
\]

\step{4}{Comparing coefficients}
The result is
\[
rX+\sum_k(-1)^kY_k\sum_{i:Z_i=Y}q_i\,\frac{\tau_i^k}{k!}=\sum_k(-1)^k\frac{t^{k+1}}{(k+1)!}\,Y_k .
\]
By independence, $r=0$ and, for every $k\ge0$,
\[
\sum_{i:Z_i=Y}q_i(t)\,\tau_i(t)^k=\frac{t^{k+1}}{k+1}.
\]

\step{5}{A rank contradiction}
Fix $t\in(0,r_0)$ and let $p=\#\{i:Z_i=Y\}$. The sequence $m_k=t^{k+1}/(k+1)=\sum_{i=1}^pq_i\tau_i^k$ has Hankel matrices
\[
H_n=[m_{i+j}]_{i,j<n}=W^{\mathsf T}\,\mathrm{diag}(q)\,W,\qquad W_{ij}=\tau_i^j,
\]
of rank at most $p$. On the other hand
\[
H_n=t\,\mathrm{diag}(t^i)\,\mathrm{Hilb}_n\,\mathrm{diag}(t^j),\qquad \mathrm{Hilb}_n=\Big[\frac1{i+j+1}\Big],
\]
and the Hilbert matrix is positive definite (the Gram matrix of monomials in $L^2(0,1)$), so $H_n$ has rank $n$. For $n=p+1$ this is a contradiction.

\whyworks{a finite product of flows can only produce finitely many ``exponential moments'', while the true first-order interaction needs infinitely many independent ones.}
\end{proof}
\uses{none.}

\begin{example}{Where the hypothesis holds, and where it does not}
In the Weyl algebra, $X=x^a\partial$ and $Y=\partial$ with $a\ge3$ give $Y_k=c_k\,x^{k(a-1)}\partial$ with $c_k\ne0$: distinct degrees, so the hypothesis holds. For $\partial^2$ and $x^k$ with $k\ge3$ it fails (each is locally nilpotent on the other), and the question is open \led{P1-Q3}.
\end{example}
